% =============================================================================
% 053 / 068
% Status: raw
% =============================================================================
% Notes:
%
% Source question:
% 吳語晚安怎麼說?
% 明朝見的意思無法和晚安對應。
% =============================================================================

\chapter[吳語晚安怎麼說? 明朝見的意思無法和晚安對應。]{吳語晚安怎麼說? \\[0.35em] 明朝見的意思無法和晚安對應。}
\qaanswer{
各地沒有統一的說法。但各地原生吳語基本不用「時間+問候語」這樣的形式。府城受漢語官話和普通話影響有用「時間+問候語」的形式，但這不是吳語本有的表現。比如北部吳語區的府城有「日安」、「夜安」的說法，屬於明清官話體系的用法。現在也逐漸被「早上好」、「夜到好」這樣的偏普通話語法的句式用法所取代。

另外「明朝 見」是官話+普通話的混合用法，府城一般用「明朝 會」，「見」、「會」都是讀原調，與「明朝」並不是黏附關係，沒有變調，意思是「再見」。
}

